\subsection{Conditional Statements}
Suppose we have in mind a specific integer a. 
Consider the following statement about a.\\

R: If the integer a is a multiple of 6, then a is divisible by 2. \\

Notice that R is built up from two simpler statements:
\begin{center}
    R: The integer a is a multiple of 6.

    Q: The integer a is divisible by 2. 

    If P, then Q
\end{center}

This is a perfect example of conditional statement. 

\begin{definition}
    [Conditional Statement]
    A statement of form $P \implies Q$ is called a conditional statement because it means Q will be true under the condition that P is true. 
\end{definition}

Think of $P \implies Q$ as a promise that whenever P is ture, Q will be true also.
There is only one way this promise can be broken,
namely if P is true but Q is false. 

There are many english sentences that are able to represent $P \implies Q$. Here is a list:
\begin{itemize}
    \item P implies Q 
    \item If P, then Q 
    \item Q if P 
    \item Q whenever P. 
    \item Whenever P, then also Q 
    \item P is a sufficient condition for Q 
    \item For Q, it is sufficient that P. 
    \item Q is a necessary condition for Q.
    \item P only if Q. 
\end{itemize}
\newpage
